% !TEX root = main.tex
% =============================================================================
% Method: hybrid trunk-share / junction-split + Hermite
% Cost formulas → sec_costs (next).
% =============================================================================

\section{Method}
\label{sec:method}

This section fixes the parameterization and the optimize loop; cost formulas
are in Sec.~\ref{sec:costs}.

\subsection{Pipeline overview}
\label{sec:method:pipeline}

For each group independently (cross-group coupling enters only through
clearance):

\begin{enumerate}
  \item \textbf{Tree.} Deduplicate the coarse grid / polyline embedding into an
        undirected tree $\mathcal{T}$ (or forest): nodes are terminals and
        branch points; edges carry the coarse polylines and radii.
  \item \textbf{Junction partition.} Grow junction complexes around
        degree-$\ge 3$ nodes; label remaining edges as trunks
        (Sec.~\ref{sec:method:partition}).
  \item \textbf{Hermite graph.} Instantiate cubic Hermite segments on labeled
        edges. Trunk knots are shared; junction knots are private per labeled
        junction edge; ports and trunk--junction interfaces are clamped for
        $C^{0}/C^{1}$ by construction (Sec.~\ref{sec:method:hermite}).
  \item \textbf{Assemble \& optimize.} Sample usages, minimize soft $J$ by
        gradient steps, accept under hard clearance (Alg.~\ref{alg:refine}).
\end{enumerate}

Groups couple only through $J_{\mathrm{clear}}$ and the final accept test.
Topology $\mathcal{T}$ and $\mathcal{U}$ are never rewritten.

\subsection{Junction partition: where DOF are private}
\label{sec:method:partition}

The hybrid rule is: \emph{share controls wherever usages share physical pipe;
add private controls only where multi-usage curvature conflicts}.

\paragraph{Grow complexes.}
Seed every tree node of degree $\ge 3$.
Grow a zone around each seed by tree-path length with threshold
$\tau=\max(\tau_0,\,\alpha\,r_{\mathrm{local}},\,\beta\,R_{\mathrm{bend}})$.
After the zone grow, any connected trunk component whose total length is shorter
than a threshold $\tau_{\mathrm{trunk}}$ is absorbed into junction (so nearby
seeds separated by a short trunk become one complex $Z$).
Edges inside $Z$ are labeled \texttt{junction}; edges outside all complexes are
\texttt{trunk}.

\paragraph{Variable layout.}
Table~\ref{tab:regions} summarizes roles after the label pass.
When building the Hermite graph we apply one privatization rule:
\emph{ports and complex interfaces keep a single global knot; every other
sample on a labeled edge is private to that edge}, so two junction edges that
meet at the same coarse branch point do not share interior $(p,v)$.
Usages do not each own a full duplicate spline; they inherit geometry by
walking the labeled edges on their tree path.
Junction chains stay short, so new freedom appears only near the conflict in
Fig.~\ref{fig:y-shape}.

\begin{table}[H]
  \centering
  \caption{Region roles after junction partition.}
  \label{tab:regions}
  \begin{tabular}{@{}p{0.16\linewidth}p{0.28\linewidth}p{0.24\linewidth}p{0.24\linewidth}@{}}
    \toprule
    Region & Geometry & Variables & Continuity \\
    \midrule
    Trunk
      & Edges outside all complexes
      & Shared Hermite knots $X_{\mathrm{trunk}}$
      & Intrinsic (one curve) \\
    Junction
      & Interior of complex $Z$
      & Private knots \emph{per labeled junction edge}
      & Meets trunk at interface \\
    Interface
      & Complex boundary
      & Shared endpoints (+ clamped tangents)
      & Hard $C^{0}$; hard $C^{1}$ by construction \\
    Ports
      & Group terminals
      & Fixed position / axis tangent
      & Hard clamps \\
    \bottomrule
  \end{tabular}
\end{table}

\paragraph{Y example (how a local split appears).}
On $A$--$J$--$C$ with $B$--$J$ and
$\mathcal{U}=\{A\!\to\!C,\,B\!\to\!C,\,A\!\to\!B\}$, the coarse embedding places
all three branches at one point $J$.
After partition, each usage is a short alternating chain of
\emph{shared trunks} and \emph{private junction edges}:
\begin{align*}
A\!\to\!C &:\;
  A_{\mathrm{trunk}}
  \;\to\;
  AC_{\mathrm{junc}}
  \;\to\;
  C_{\mathrm{trunk}},
\\
B\!\to\!C &:\;
  B_{\mathrm{trunk}}
  \;\to\;
  BC_{\mathrm{junc}}
  \;\to\;
  C_{\mathrm{trunk}},
\\
A\!\to\!B &:\;
  A_{\mathrm{trunk}}
  \;\to\;
  AB_{\mathrm{junc}}
  \;\to\;
  B_{\mathrm{trunk}}.
\end{align*}
Here $A_{\mathrm{trunk}}$, $B_{\mathrm{trunk}}$, and $C_{\mathrm{trunk}}$ are
shared Hermite geometry (every usage that enters that arm reads the same knots).
The three junction edges $AC_{\mathrm{junc}}$, $BC_{\mathrm{junc}}$, and
$AB_{\mathrm{junc}}$ are independent: each has its own interior $(p,v)$, even
though they all start at the same coarse coordinates near $J$.
Thus $A\!\to\!C$ and $B\!\to\!C$ couple through $C_{\mathrm{trunk}}$, while
$A\!\to\!B$ can fair through $AB_{\mathrm{junc}}$ without forcing a hairpin at a
single shared point.
Initialization copies the shared-junction layout, so the private junction knots
begin coincident; during optimization they may separate and open the void in
Fig.~\ref{fig:y-shape}b---a continuous motion of already-private DOF, not a
discrete mode switch.

\subsection{Cubic Hermite representation}
\label{sec:method:hermite}

Each labeled edge is a chain of cubic Hermite segments.
A segment with ends $(p_0,v_0)$ and $(p_1,v_1)$ is, for $t\in[0,1]$,
\begin{equation}
\label{eq:hermite}
\gamma(t)
=
H_{00}(t)\,p_0
+ H_{10}(t)\,v_0
+ H_{01}(t)\,p_1
+ H_{11}(t)\,v_1,
\end{equation}
with the standard Hermite basis
\begin{align*}
H_{00}(t) &= 2t^{3}-3t^{2}+1, &
H_{10}(t) &= t^{3}-2t^{2}+t, \\
H_{01}(t) &= -2t^{3}+3t^{2}, &
H_{11}(t) &= t^{3}-t^{2}.
\end{align*}
Thus $\gamma(0)=p_0$, $\gamma(1)=p_1$, $\gamma'(0)=v_0$, $\gamma'(1)=v_1$.
Shared nodes enforce hard $C^{0}$ and $C^{1}$: every incident segment reads the
same $(p,v)$.

\paragraph{Clamps.}
Port nodes fix $p$ and align $v$ with the prescribed unit vector.
Interface nodes between trunk and junction share $p$ (and, by construction,
compatible $v$) so that usage walks stay $G^{1}$ across the complex boundary.
Only movable trunk knots and private junction knots enter the decision
vector $x$.

\subsection{Assemble and optimize}
\label{sec:method:opt}

Alg.~\ref{alg:refine} minimizes the soft objective of Sec.~\ref{sec:costs} by
first-order steps on $\nabla_x J$ and accepts a design only under hard obstacle /
cross-group clearance.
Cost formulas are in Sec.~\ref{sec:costs}; which first-order variant we use in
practice is an implementation detail reported with the experiments.

\begin{algorithm}[H]
\caption{Usage assembly and gradient refine (per multi-group instance)}
\label{alg:refine}
\begin{algorithmic}[1]
\Require Hermite graphs $\{G_g\}$, usages $\mathcal{U}$, weights, steps $T$, refresh period $k$
\State $x \leftarrow$ concatenate free node positions and tangents $(p,v)$ across groups
\For{$t = 1,\ldots,T$}
  \State write $x$ into the Hermite graphs
  \If{$t=1$ or $(t-1) \bmod k = 0$}
    \State find candidate colliding pairs by proximity (broadphase):
    \Statex \quad $\mathcal{N}_{\mathrm{obs}}$: tube samples vs obstacle primitives
    \Statex \quad $\mathcal{N}_{\mathrm{grp}}$: tube samples of group $g$ vs group $g'$
    \Statex \quad $\mathcal{N}_{\mathrm{union}}$: nearby junction samples within a complex
  \EndIf
  \For{each usage $u \in \mathcal{U}$}
    \State follow $u$'s path on $\mathcal{T}$ (shared trunks + private junctions)
    \State sample Hermite segments along that path $\to \gamma_u,\, r_u$
  \EndFor
  \State $J_{\mathrm{path}} \leftarrow$ length / curvature on $\{\gamma_u,r_u\}$
        \Comment{Sec.~\ref{sec:costs}}
  \State $J_{\mathrm{space}} \leftarrow$ spacing penalty on chord gaps
  \State $J_{\mathrm{clear}} \leftarrow$ soft gaps on pairs in
        $\mathcal{N}_{\mathrm{obs}}\cup\mathcal{N}_{\mathrm{grp}}$
  \State $J_{\mathrm{union}} \leftarrow$ soft compactness on pairs in
        $\mathcal{N}_{\mathrm{union}}$
  \State $J \leftarrow J_{\mathrm{path}}+J_{\mathrm{space}}+J_{\mathrm{clear}}+J_{\mathrm{union}}$
  \State differentiate $J$ w.r.t.\ $x$ through the sampling map (reverse-mode AD)
  \State $x \leftarrow$ gradient-based update using $\nabla_x J$
\EndFor
\State write $x$ back; hard-check obstacle / cross-group clearance (no soft margin)
\State \textbf{accept} iff hard hit count $= 0$; else reject / raise clearance weights
\Ensure accepted Hermite graphs, or failure
\end{algorithmic}
\end{algorithm}
